\paragraph{From the discriminant to a height.}
\hypertarget{source-slice-heights}{}
We give the height calculation before the final lift. Use block
sphere probability measure on the rows of complex coefficients
and the Gauss norm at the finite places. For a polynomial \(q\)
let \(h_0(q)\) denote the sum of its finite log Gauss norms and
the integrals of \(\log|q|\) on the product of spheres, with the
number-field normalization on both sums. On constants the sum
is zero. The height is additive on products, so extends to
rational functions. Let \(N\) be the number of light columns
and \(\ell_N=\log N+1/2\).

\begin{lemma}[The coefficient-sphere height]
\label{lem:coefficient-sphere-height}
If \(q\ne0\) has degrees \(m_i\)
homogeneously in the respective rows, then
\begin{equation}\label{eq:slice-sphere-lower}
 h_0(q)\ge-\ell_N\sum_i m_i .
\end{equation}
\end{lemma}

\begin{proof}
Indeed fix a lexicographic order on its monomials. Condition
at each sphere on the absolute values of its coordinates.
Their phases are independent uniform angles. Successive use
of Jensen's inequality for a univariate polynomial on the
circle bounds the phase integral below by the logarithm of
the absolute value of the lexicographically leading term.
For a uniform complex unit vector the mean log absolute
value of one coordinate is \(-H_{N-1}/2\).

For \(N>1\)
its squared absolute value has density
\((N-1)(1-t)^{N-2}\) on \((0,1)\). Integrating \(\log t\)
gives this formula, including the value zero when \(N=1\).
For \(N>1\) the one-coordinate expectation is
\[
  \mathbb E\log|z_1|=\frac{N-1}{2}\int_0^1(1-t)^{N-2}\log t\,dt=-\frac{H_{N-1}}2 .
\]
The same leading coefficient is available in every Gauss
norm, so the product formula cancels its logarithms.
Finally \(H_{N-1}/2\le\ell_N\). The conditional Jensen
calculation proves integrability because the positive part on
the spheres is bounded and the lower integral is finite.

\end{proof}

\begin{proposition}[The discriminant of the selected slice]
\label{prop:slice-height-discriminant}
For the selected finite reduced slice, the weighted discriminant term satisfies \eqref{eq:slice-height-disc},
\[
 2\langle I\rangle\le\delta\log\delta+2\sum_i H(t'_i)+2s a_F\ell_N+2H_F.
\]
 An integral auxiliary section
on those same branches contributes through its normalized height and
its mean logarithmic supremum, with the same weight line and branch
multiplicities.
\end{proposition}

\begin{proof}
\proofstep{The discriminant in a trivialization of the weight line.}
For a section \(\tau\) nonzero on each component let
\(n_q(\tau)\) be the sum of its integral frame orders over
places of the incidence fields above the divisor \(q=0\)
of \(K[X]\), with residue degrees. A place trivial on the
component field contributes zero in the constant generic
frame. Define \(J_\tau=\prod_q q^{n_q(\tau)}\), and
\(a_\tau=\sum_q n_q(\tau)\deg_1 q\) where row 1 is one
of the \(s\) cut rows. These horizontal divisors have finite support, are independent
of row zero and have equal degree in each cut row for each
contributing irreducible \(q\).

We prove these properties
below by extension of the incidence point to its divisorial
valuation rings and symmetry of the cut. The height of a framed
weight is
\[
 H_\tau=\langle N(\tau)\rangle
       +\operatorname{mean}_{\sigma}\int
            \sum_{P}-\log\|\tau(P)\|_\sigma
       +h_0(J_\tau),\qquad
 N_\lambda(\tau)=\sum_{Z,w\mid\lambda} f_w n_w(\tau).
\]
Here \(P\) runs over complex points of the generic slice
at the coefficient tuple being integrated. Discriminants
and exceptional denominator loci are null sets. The integrability argument below applies to each of the
metric terms on this family.
Changing the frame within a weight adds the orders and
log norms of the field norm of its componentwise ratio.
The product identity for that rational function cancels
the change in the three terms. Likewise heights and
horizontal degrees add upon multiplication of weights
and sections. \label{cor:slice-heights-1}

\proofstep{Horizontal orders and the incidence divisor.}
Form 
\[
  U=\Delta J_\tau^2.
\]
 At a horizontal prime apply
the trace calculation of \eqref{eq:slice-discriminant-order}
in the corresponding DVR. At the contributing places the
short sections have nonnegative orders in the integral
frame. At the other places a nonzero constant component
reading has order zero. Hence the local order of \(U\)
is nonnegative at every irreducible of \(K[X]\). Unique
factorization makes \(U\) a polynomial.
In reduced form \(U=A/B\) any irreducible \(q\mid B\) would have
\[
  \operatorname{ord}_q U=\operatorname{ord}_q A-\operatorname{ord}_q B=-\operatorname{ord}_q B<0 .
\]
 Its row degrees
are \((0,2a_\tau,\ldots,2a_\tau)\) because the equal-weight
coordinates of \(b_i\) depend only on the cut points, so \(\Delta\) is invariant under
independent scalar changes of the rows. For example the
corresponding substitutions on the incidence fields fix the
\(b_i\) and intertwine the traces. All the degree therefore
comes from the horizontal factor.

\proofstep{The trace determinant in the archimedean average.}
For almost every tuple the trace discriminant is the square
of the determinant of the matrix \(b_i(P)\). Multiply each
row by the reading of \(\tau(P)\) in a unit norm frame.
Hadamard on columns bounds the squared absolute determinant
of this new matrix by \(\delta^\delta
\prod_i\|t'_i\|_{\sup,\sigma}^2\).
In the finite height of \(\Delta\) the log Gauss norm has the
opposite sign to \(v_\lambda(\Delta)\log\mathrm N\lambda\).
Additivity of \(h_0\) gives the exact identity
\[
 \langle v(\Delta)+2N(\tau)\rangle
 = 2H_\tau-h_0(U)
  +\operatorname{mean}_{\sigma}\int
       \left(\log|\Delta|+2\sum_P\log\|\tau(P)\|_\sigma\right).
\]
Combining the preceding inequalities yields
\begin{equation}\label{eq:slice-height-disc}
 2\langle I\rangle\le \delta\log\delta+2\sum_i H(t'_i)
                     +2s a_F\ell_N+2H_F .
\end{equation}
Here the subscript \(F\) means the weight \(w_F\).
For an integral section \(g\) with nonnegative horizontal
orders and bounded norm the same height formula gives
\[
 0\le\langle N(g)\rangle\le H_g+
    \delta\operatorname{mean}_\sigma\log\|g\|_{\sup,\sigma}
                       +s a_g\ell_N .
\]
Apply this to \(\eta\) and divide by \(i_\eta\).
Thus both vertical prices in \eqref{eq:geometric-four-terms}
are bounded by heights of the same global family.

\end{proof}

\paragraph{The degree and the auxiliary height.}

The preceding height contains contributions of different degrees. Divide by the incidence degree before estimating the auxiliary term.

For the following comparisons set \(n_1=i_\eta g_0^*\).
Write \(H_e=H_\eta/n_1\), \(a_e=a_\eta/n_1\).
Use the integral sections \(t_1,t_2\) of weights
\(2(w_L-w_\eta/n_1)\) and \((s+4)w_L-w_F\).
\(H_1,a_1\) denote the corresponding quantities for
\(t_1\) divided by two. Choose a constant light combination
nonzero on every component. Such combinations form the
complement of finitely many proper subspaces over an infinite
field. Let \(a_L,H_L\) be its degree and height.
By the weight identities,
\[
 a_L=a_1+a_e,\qquad (s+4)a_L=a_F+a_2 .
\]
All degrees on the right are nonnegative by the integral
section comparison. For \(a_F\) one can take a short
section nonzero on each particular component and use frame
independence.

The normalized moving part of the light
system has degree \(\delta\) because adding a general zeroth
hyperplane to the \(s\) cutting hyperplanes yields the
incidence form of the light image, with its generic field
multiplicity on each component. Equivalently the norm of
the ratio of two such hyperplanes is the ratio of their
products over the \(\delta\) slice points. After removing
the minimum order of the light columns at each place, its
specialization gives the divisor of that moving system.
The actual degree is 
\[
  a_L=\delta+\sum_q E_q\deg_1 q,
\]
where \(E_q\) is the sum of the removed minima with residue
degrees. On the horizontal depth estimate
\(0\le E_q\le c\,n_q(\eta)/(2i_\eta)\) we get
\[
 a_L\le\delta+\tfrac12 c g_0^* a_e
      \le\delta+\tfrac14 a_L .
\]
In particular \(a_L\le4\delta/3\) and
\(a_F\le4(s+4)\delta/3\).

\begin{proposition}[Heights of the chosen incidence sections]
\label{prop:slice-normalized-heights}
Under the degree, metric and Chow bounds used below, for \(s\ge3\) and
\(D_0=\delta S_L\) the pre-final heights satisfy
\[
 H_e\le(s+1)D_0(2G_c+6P_H),\qquad
 H_F\le(s+4)(s+1)D_0(2G_c+6P_H).
\]
The definite lift satisfies the same upper bound for \(H_F\), using
the baseline finite depths alone.
\end{proposition}

\begin{proof}
\proofstep{The normalized height of the moving incidence divisor.}
For precision, let \(H_{\rm mov}\) be the height of the
moving light system, computed at finite places by subtracting the minimum of the
light orders, and at a complex point divide light norms by
\((\sum_l\|w_l\|^2)^{1/2}\). Its incidence height expression
is
\[
 H_{\rm mov}=\sum_Z h_0(F_Z)+\delta H_{N-1}/2 ,
\]
where \(F_Z\) is the Chow form of the light image raised to
the generic covering degree. In fact the zeroth row in
\(F_Z\), with the other rows generic, factors up to scalar
as the product of the light linear forms at the slice
points.

At a Gauss place each factor has the minimum
column order. The mean log absolute value of each factor
on the complex unit sphere is the log norm of its vector
minus \(H_{N-1}/2\), by unitary invariance.
For a nonzero coefficient vector \(v\), the sphere average is therefore
\[
  \int_{\|z\|=1}\log|\langle z,v\rangle|\,d\sigma(z)=\log\|v\|-\tfrac12 H_{N-1}.
\]
 Comparing
with a fixed zeroth row at which the product is nonzero
gives the displayed formula by additivity of the sphere
and Gauss heights. This calculation uses the specialization
of the horizontal divisor at that row.

\proofstep{Using the same integral sections in the height comparison.}
To get the height bounds, form \(H_{\rm mov}\) using exactly the
light-weight columns retained in the new choice of sections. Write \(L=S_L\) and
\(\beta=cg_0^*\le1/2\).
Let \(P_0,P_N,P_5,P_2,P_4\ge1\) bound respectively
the integral-section price in the Chow estimate, the
finite depths, the two terms comparing the archimedean light
metrics, and the norms of the two complementary sections.
Put 
\[
  p=P_0+P_N+P_5+P_2+P_4+P_{rw},
\]
and take \(P_H\ge 2p+G_c\). Use the following inputs in the normalized height calculation.\begin{align*}
 \sum_Z h_0(F_Z)&\le (s+2)\delta
              (P_0L+G_cL+\log N+1),\\
 H_L&\le H_{\rm mov}
    +\delta(P_N+P_5+P_2)L
    +s\beta a_e\ell_N+\beta H_e,\\
 H_e&\le H_L+\delta P_4s L+s a_1\ell_N,\\
 H_F&\le(s+4)H_L+\delta P_4s L+s a_2\ell_N .
\end{align*}

\proofstep{The finite contribution of the auxiliary section.}
For the second inequality, let \(B_\lambda\) be the sum of the
minimum light orders with residue degrees at a finite place
\(\lambda\). At each branch its baseline is
\(e_w\nu_\lambda\). The decomposition identity
\(\sum_{Z,w\mid\lambda}e_w f_w=\delta\) gives
\[
 B_\lambda\le\delta\nu_\lambda+
                   \frac{c}{2i_\eta}N_\lambda(\eta).
\]
Use the weighted baseline bound
\(\langle\nu\rangle\le P_N L\) on the run. Since the
auxiliary orders are nonnegative, one can use the larger
coefficient \(c/i_\eta\) for this finite contribution.
At a complex slice point the change from the moving metric
contributes \(-\log(\sum_l\|w_l\|^2)^{1/2}\). The paired light
and auxiliary metric bound pays \(P_5D_F+cP_2\) and the same
coefficient of \(-\log\|\eta\|\).

Horizontally, let
\(J_{\min}=\prod_q q^{E_q}\) in the frame of the light system.
The effective difference with coefficients
\((c/i_\eta)n_q(\eta)-E_q\) has sphere height at least
\(-s(c/i_\eta)a_\eta\ell_N\), by
\eqref{eq:slice-sphere-lower}. This also holds for rational
or real exponents by summing the individual bounds with
nonnegative multipliers. Therefore
\[
 h_0(J_{\min})\le (c/i_\eta)h_0(J_\eta)+s\beta a_e\ell_N .
\]
In the three terms of \(H_L-H_{\rm mov}\) the resulting
auxiliary coefficient is exactly \(c/i_\eta=\beta/n_1\).
The last two height inequalities follow by weight additivity
and the lower bound for the height of an integral section.
Use the bound \(P_4sL\) for the two log suprema in
Paragraph~\ref{par:selected-height-sections}, with the factor
one half for \(t_1\).

\proofstep{Bounding the coefficient-sphere integrals.}
For example put \(D_0=\delta L\) in this calculation.
Use \(\log N+1,\ell_N,H_{N-1}/2\le P_{rw}L\).
The first two inequalities imply
\[
 H_L\le D_0\{(s+2)(p+G_c)+p
                   +\tfrac23sP_{rw}\}+\beta H_e .
\]
Indeed the harmonic and depth terms not in the first
parenthesis are paid by \(p\). The light-count term in
the first parenthesis leaves further slack. Since
\(H_e\le H_L+D_0 s(P_4+\tfrac43P_{rw})\),
substitution and division by \(1-\beta\) give
\[
 H_L\le D_0\{2(s+2)(p+G_c)+2p+
                            s(P_4+\tfrac83P_{rw})\}.
\]
If \(H_L\) is negative the same upper bound follows
directly. Adding \(D_0s(P_4+\tfrac43P_{rw})\)
costs at most \(D_0s(2P_4+4P_{rw})\) in total
in the last \(s\)-term. Hence, for \(s\ge3\),
\[
 H_e\le(s+1)D_0(2G_c+6P_H),\qquad
 H_F\le(s+4)(s+1)D_0(2G_c+6P_H).
\]
The second inequality uses \(a_2\le (s+4)4\delta/3\).

At the definite final lift use the baseline depth bound alone.
The light minimum has zero horizontal contribution, so
\(a_L=\delta\) and \(a_F\le(s+4)\delta\). The first-minimum
function is zero and the lower bound for the light vector is
uniform over the complex frames. In the height comparison this
gives
\[
 H_L\le H_{\rm mov}+\delta(P_N+P_5)S_L .
\]
The complementary section \(t_2\) gives
\(H_F\le(s+4)H_L+\delta P_4sS_L+s a_2\ell_N\). With \(a_2\le(s+4)\delta\), the Chow bound and
\(\ell_N\le P_{rw}S_L\) give
\[
 H_F\le(s+4)(s+1)\delta S_L(2G_c+6P_H).
\]
Indeed \(H_{\rm mov}\) contributes at most
\(\delta S_L((s+2)(P_0+G_c+P_{rw})+P_{rw})\).

The other terms after dividing by \((s+4)\delta S_L\)
are at most \(P_N+P_5+P_4s/(s+4)+sP_{rw}\).
Their sum is within the stated bound by \(P_H\ge2p+G_c\).
The light and complementary section norms at this lift are
integrable using the same rational ratios and the bounded light
vector above and below. Thus the discriminant argument uses
the same upper bound at \(j=m\) with no auxiliary order term.
\label{cor:slice-heights-2}

Together with \eqref{eq:slice-height-disc} and

\[
  \operatorname{mean}\log\|\eta\|_{\sup}/i_\eta\le P_\eta
\]
these are the two slice-height estimates displayed above.
We next derive the Chow bound and transfer the metric estimates
to the places used here. We then justify horizontal specialization and
integrability on this incidence correspondence.
\label{cor:slice-heights-3}
\end{proof}
